% GENERATED FILE. Edit canonical lessons, then run npm run generate:books.
% canonical-source-hash: fnv1a64:cd046fc52bd28fa1
% canonical-lessons: JA-C443-tekubi, JA-C443-ashikubi, JA-C443-oyayubi, JA-C443-kuchibiru, JA-C443-mabuta, JA-R443-review-438-443

\chapter{Wrists, Ankles and Lips}
\label{ch:ja-wrists-ankles-and-lips}
\hlchaptermodality{443}

\begin{chapteropening}
\textbf{By the end of this chapter:} \emph{I can name small parts of the body.}

Complete the last lesson of chapter 443: I can name small parts of the body.
\end{chapteropening}

\section[\texorpdfstring{\ja{てくび} (tekubi)}{てくび (tekubi)}]{\ja{てくび} (tekubi) — a wrist}

\label{lesson:JA-C443-tekubi}



\begin{quote}
Before the new one: say the Japanese for the spine, then the Japanese for a joint.
\end{quote}

\subsection*{You'll want to know: \ja{てくび}}
\textbf{\ja{てくび}} — \emph{tekubi} — \textquotedblleft{}a wrist\textquotedblright{}.

\begin{quote}
\textbf{\ja{ころんで、てくびが} \ja{いたいです。}} — \emph{koronde, tekubi ga itai desu} — \textquotedblleft{}I fell over, and my wrist hurts.\textquotedblright{}
\end{quote}

\begin{cousinweb}
\textbf{\ja{てくび}} (\emph{tekubi}) is a native Japanese compound (wago), inherited rather than borrowed from Chinese: \emph{te}, \textquotedblleft{}hand\textquotedblright{} and \emph{kubi}, \textquotedblleft{}neck\textquotedblright{} (JMdict, KANJIDIC2).
\end{cousinweb}

\begin{grammarlens}[title={What you've built}]
One more word for describing what you lived and felt.
\end{grammarlens}

\subsection*{Your turn}
\begin{itemize}
\raggedright
  \item \emph{Say it:} \emph{tekubi}
  \item \emph{Say it:} \emph{tekubi}, once more
  \item \emph{Say it:} \emph{koronde, tekubi ga itai desu} — I fell over, and my wrist hurts.
  \item \emph{Recall:} say the Japanese for the spine, then the Japanese for a joint, then say \emph{tekubi} again
\end{itemize}

\begin{tcolorbox}[breakable,colback=teal!4,colframe=teal!35!black,title={Before you move on}]
What does \ja{てくび} mean? (\textquotedblleft{}A wrist\textquotedblright{}.) Say it once more.
\end{tcolorbox}
\section[\texorpdfstring{\ja{あしくび} (ashikubi)}{あしくび (ashikubi)}]{\ja{あしくび} (ashikubi) — an ankle}

\label{lesson:JA-C443-ashikubi}



\begin{quote}
Before the new one: say the Japanese for a joint, then the Japanese for a wrist.
\end{quote}

\subsection*{You'll want to know: \ja{あしくび}}
\textbf{\ja{あしくび}} — \emph{ashikubi} — \textquotedblleft{}an ankle\textquotedblright{}.

\begin{quote}
\textbf{\ja{はしって、あしくびが} \ja{いたく} \ja{なりました。}} — \emph{hashitte, ashikubi ga itaku narimashita} — \textquotedblleft{}I ran, and my ankle started to hurt.\textquotedblright{}
\end{quote}

\begin{cousinweb}
\textbf{\ja{あしくび}} (\emph{ashikubi}) is a native Japanese compound (wago), inherited rather than borrowed from Chinese: \emph{ashi}, \textquotedblleft{}foot\textquotedblright{} and \emph{kubi}, \textquotedblleft{}neck\textquotedblright{} (JMdict, KANJIDIC2).
\end{cousinweb}

\begin{grammarlens}[title={What you've built}]
One more word for describing what you lived and felt.
\end{grammarlens}

\subsection*{Your turn}
\begin{itemize}
\raggedright
  \item \emph{Say it:} \emph{ashikubi}
  \item \emph{Say it:} \emph{ashikubi}, once more
  \item \emph{Say it:} \emph{hashitte, ashikubi ga itaku narimashita} — I ran, and my ankle started to hurt.
  \item \emph{Recall:} say the Japanese for a joint, then the Japanese for a wrist, then say \emph{ashikubi} again
\end{itemize}

\begin{tcolorbox}[breakable,colback=teal!4,colframe=teal!35!black,title={Before you move on}]
What does \ja{あしくび} mean? (\textquotedblleft{}An ankle\textquotedblright{}.) Say it once more.
\end{tcolorbox}
\section[\texorpdfstring{\ja{おやゆび} (oyayubi)}{おやゆび (oyayubi)}]{\ja{おやゆび} (oyayubi) — a thumb}

\label{lesson:JA-C443-oyayubi}



\begin{quote}
Before the new one: say the Japanese for a wrist, then the Japanese for an ankle.
\end{quote}

\subsection*{You'll want to know: \ja{おやゆび}}
\textbf{\ja{おやゆび}} — \emph{oyayubi} — \textquotedblleft{}a thumb\textquotedblright{}.

\begin{quote}
\textbf{\ja{おやゆびに} \ja{きずが} \ja{あります。}} — \emph{oyayubi ni kizu ga arimasu} — \textquotedblleft{}There is a cut on my thumb.\textquotedblright{}
\end{quote}

\begin{cousinweb}
\textbf{\ja{おやゆび}} (\emph{oyayubi}) is a native Japanese compound (wago), inherited rather than borrowed from Chinese: \emph{oya}, \textquotedblleft{}parent\textquotedblright{} and \emph{yubi}, \textquotedblleft{}finger\textquotedblright{} (JMdict, KANJIDIC2).
\end{cousinweb}

\begin{grammarlens}[title={What you've built}]
One more word for describing what you lived and felt.
\end{grammarlens}

\subsection*{Your turn}
\begin{itemize}
\raggedright
  \item \emph{Say it:} \emph{oyayubi}
  \item \emph{Say it:} \emph{oyayubi}, once more
  \item \emph{Say it:} \emph{oyayubi ni kizu ga arimasu} — There is a cut on my thumb.
  \item \emph{Recall:} say the Japanese for a wrist, then the Japanese for an ankle, then say \emph{oyayubi} again
\end{itemize}

\begin{tcolorbox}[breakable,colback=teal!4,colframe=teal!35!black,title={Before you move on}]
What does \ja{おやゆび} mean? (\textquotedblleft{}A thumb\textquotedblright{}.) Say it once more.
\end{tcolorbox}
\section[\texorpdfstring{\ja{くちびる} (kuchibiru)}{くちびる (kuchibiru)}]{\ja{くちびる} (kuchibiru) — lips; a lip}

\label{lesson:JA-C443-kuchibiru}



\begin{quote}
Before the new one: say the Japanese for an ankle, then the Japanese for a thumb.
\end{quote}

\subsection*{You'll want to know: \ja{くちびる}}
\textbf{\ja{くちびる}} — \emph{kuchibiru} — \textquotedblleft{}lips; a lip\textquotedblright{}.

\begin{quote}
\textbf{\ja{ふゆは} \ja{くちびるが} \ja{かわきます。}} — \emph{fuyu wa kuchibiru ga kawakimasu} — \textquotedblleft{}My lips get dry in winter.\textquotedblright{}
\end{quote}

\begin{cousinweb}
\textbf{\ja{くちびる}} (\emph{kuchibiru}) is a native Japanese word (wago), inherited rather than borrowed from Chinese; its kanji spelling was attached to it later (JMdict).
\end{cousinweb}

\begin{grammarlens}[title={What you've built}]
One more word for describing what you lived and felt.
\end{grammarlens}

\subsection*{Your turn}
\begin{itemize}
\raggedright
  \item \emph{Say it:} \emph{kuchibiru}
  \item \emph{Say it:} \emph{kuchibiru}, once more
  \item \emph{Say it:} \emph{fuyu wa kuchibiru ga kawakimasu} — My lips get dry in winter.
  \item \emph{Recall:} say the Japanese for an ankle, then the Japanese for a thumb, then say \emph{kuchibiru} again
\end{itemize}

\begin{tcolorbox}[breakable,colback=teal!4,colframe=teal!35!black,title={Before you move on}]
What does \ja{くちびる} mean? (\textquotedblleft{}Lips; a lip\textquotedblright{}.) Say it once more.
\end{tcolorbox}
\section[\texorpdfstring{\ja{まぶた} (mabuta)}{まぶた (mabuta)}]{\ja{まぶた} (mabuta) — an eyelid}

\label{lesson:JA-C443-mabuta}



\begin{quote}
Before the new one: say the Japanese for a thumb, then the Japanese for lips.
\end{quote}

\subsection*{You'll want to know: \ja{まぶた}}
\textbf{\ja{まぶた}} — \emph{mabuta} — \textquotedblleft{}an eyelid\textquotedblright{}.

\begin{quote}
\textbf{\ja{ねむいので、まぶたが} \ja{おもいです。}} — \emph{nemui node, mabuta ga omoi desu} — \textquotedblleft{}I am sleepy, so my eyelids are heavy.\textquotedblright{}
\end{quote}

\begin{cousinweb}
\textbf{\ja{まぶた}} (\emph{mabuta}) is a native Japanese word (wago), inherited rather than borrowed from Chinese; its kanji spelling was attached to it later (JMdict).
\end{cousinweb}

\begin{grammarlens}[title={What you've built}]
One more word for describing what you lived and felt.
\end{grammarlens}

\subsection*{Your turn}
\begin{itemize}
\raggedright
  \item \emph{Say it:} \emph{mabuta}
  \item \emph{Say it:} \emph{mabuta}, once more
  \item \emph{Say it:} \emph{nemui node, mabuta ga omoi desu} — I am sleepy, so my eyelids are heavy.
  \item \emph{Recall:} say the Japanese for a thumb, then the Japanese for lips, then say \emph{mabuta} again
\end{itemize}

\begin{tcolorbox}[breakable,colback=teal!4,colframe=teal!35!black,title={Before you move on}]
What does \ja{まぶた} mean? (\textquotedblleft{}An eyelid\textquotedblright{}.) Say it once more.
\end{tcolorbox}
\section[(dialogue)]{Looking back: from Stubborn and Hasty to Wrists, Ankles and Lips}

\label{lesson:JA-R443-review-438-443}



\begin{quote}
Say the Japanese for stubborn and impatient.
\end{quote}

\subsection*{Your turn}
Say these aloud:

\begin{itemize}
\raggedright
  \item stubborn, impatient, short-tempered, leisurely, to act big
  \item timid, a weakling, unfair, rough, courageous
  \item cautious, thoughtlessly, sloppy, neatly, suitable
  \item to boast, cheeky, admirable, a lazy person, to fool around
  \item a blood vessel, a muscle, blood, the spine, a joint
  \item a wrist, an ankle, a thumb, lips, an eyelid
\end{itemize}

\begin{tcolorbox}[breakable,colback=teal!4,colframe=teal!35!black,title={Before you move on}]
Stubborn? (\textbf{\ja{がんこ}}, \emph{ganko}.) Impatient? (\textbf{\ja{せっかち}}, \emph{sekkachi}.) Short-tempered? (\textbf{\ja{たんき}}, \emph{tanki}.) Leisurely? (\textbf{\ja{のんびり}}, \emph{nonbiri}.) To act big? (\textbf{\ja{いばる}}, \emph{ibaru}.) Timid? (\textbf{\ja{おくびょう}}, \emph{okubyō}.) A weakling? (\textbf{\ja{よわむし}}, \emph{yowamushi}.) Unfair? (\textbf{\ja{ひきょう}}, \emph{hikyō}.) Rough? (\textbf{\ja{らんぼう}}, \emph{ranbō}.) Courageous? (\textbf{\ja{ゆうかん}}, \emph{yūkan}.) Cautious? (\textbf{\ja{しんちょう}}, \emph{shinchō}.) Thoughtlessly? (\textbf{\ja{うっかり}}, \emph{ukkari}.) Sloppy? (\textbf{\ja{だらしない}}, \emph{darashinai}.) Neatly? (\textbf{\ja{きちんと}}, \emph{kichinto}.) Suitable? (\textbf{\ja{てきとう}}, \emph{tekitō}.) To boast? (\textbf{\ja{じまんする}}, \emph{jiman suru}.) Cheeky? (\textbf{\ja{なまいき}}, \emph{namaiki}.) Admirable? (\textbf{\ja{えらい}}, \emph{erai}.) A lazy person? (\textbf{\ja{なまけもの}}, \emph{namakemono}.) To fool around? (\textbf{\ja{ふざける}}, \emph{fuzakeru}.) A blood vessel? (\textbf{\ja{けっかん}}, \emph{kekkan}.) A muscle? (\textbf{\ja{きんにく}}, \emph{kinniku}.) Blood? (\textbf{\ja{けつえき}}, \emph{ketsueki}.) The spine? (\textbf{\ja{せぼね}}, \emph{sebone}.) A joint? (\textbf{\ja{かんせつ}}, \emph{kansetsu}.) A wrist? (\textbf{\ja{てくび}}, \emph{tekubi}.) An ankle? (\textbf{\ja{あしくび}}, \emph{ashikubi}.) A thumb? (\textbf{\ja{おやゆび}}, \emph{oyayubi}.) Lips? (\textbf{\ja{くちびる}}, \emph{kuchibiru}.) An eyelid? (\textbf{\ja{まぶた}}, \emph{mabuta}.)
\end{tcolorbox}
